\section{Experimental design} \label{sec:exp}
This section presents the setting of our experimental setup, in order to provide reproducibility of our results. The full data collection is made available in Github \prov{XXX}. \TODO{public repository URL plus an archived DOI (Zenodo); TR-E increasingly expects a citable archive rather than a bare link.}

\subsection{Nominal parameters}
All experiments share the vehicle and operational parameters listed in Table~\ref{tab:params}.

\begin{table}[ht]
\centering
\caption{Nominal vehicle and operational parameters.}
\label{tab:params}
\begin{tabular}{llr}
\toprule
\textbf{Category} & \textbf{Parameter} & \textbf{Value} \\
\midrule
\multirow{4}{*}{Vehicle}
  & Battery capacity $E^{\mathrm{cap}}$        & 500\,kWh \\
  & Minimum SoC $E^{\min}$                      & $0.2\cdot E^{\mathrm{cap}}$ (100\,kWh) \\
  & Full-charge time (0\%--100\%) $\mathcal{T}_R$ & 2.5\,h \\
  & \NEW{Consumption at nominal speed $\mathrm{ECR}(80)$} & \NEW{1.10\,kWh/km} \\
\midrule
\multirow{3}{*}{Charging}
  & PWL charging segments                       & 3 \\
  & PWL breakpoints (SoC)                       & 0\%, 40\%, 80\%, 100\% \\
  & \NEW{Minimum charging time $\mathcal{T}_{\min}$} & \NEW{15\,min} \\
\midrule
\multirow{6}{*}{Operations}
  & Nominal average speed                       & 80\,km/h \\
  & Customer service time                       & 30\,min \\
  & Maneuver time when seq. operations at CS stop               & 5\,min \\
  & Maneuver time at CS and rest areas stop               & 10\,min \\
  & Out of window penalty $\beta$ & 1\,h \\
  & \NEW{Idle waiting} & \NEW{not permitted} \\
\midrule
\multirow{3}{*}{Infrastructure}
  & Charging station spacing                    & \DEL{$U[40, 80]$\,km} \FIX{60\,km grid, jittered $U[-19,+19]$\,km} \\
  & Charging station or rest areas spacing      & $<$25\,km \\
  & Queue delay at CS (log-normal)              & $\mu = 10$\,min, $\sigma = 8$\,min \\
\bottomrule
\end{tabular}
\end{table}

\begin{acode}
\textbf{Charging station spacing.} \texttt{instances.py} places stations on a
deterministic 60\,km grid and jitters each position by $U[-19,+19]$\,km, so
realized spacings run roughly $[22,98]$\,km with mean 60. $U[40,80]$ was never
what the generator did, and it contradicts Section~\ref{sec:instances}'s own
``regular average intervals of 60\,km''.
\end{acode}

Charging station spacing and charger power output is taken from the Regulation (EU) 2023/1804, which mandates heavy-duty vehicles charging station every 60 km on the TEN-T core network by 2030. The energy consumption rate (ECR) as a function of speed $v$ (km/h) is modeled from \citet{gao_estimating_2025} and calibrated with verified consumption values from constructors.

\begin{acode}
\textbf{Which ECR reference?} The paper cites \citet{gao_estimating_2025};
\texttt{src/settings.py} cites Younes et al.\ (2020),
\texttt{10.1016/j.est.2020.101758}, for the functional form
$\mathrm{ECR}(v) = A/v + B + Cv^2$, cross-checked against Nykvist \& Olsson
(2021) and NACFE Run-on-Less. These cannot both be the source of the
coefficients. \TODO{resolve, and give $A$, $B$, $C$ explicitly --- the entire
energy dimension rests on three numbers that appear nowhere in the paper.}
\end{acode}

Travel time on each leg is uncertain through a multiplicative multiplier $\tilde D_i = D_i\,\xi_i$. The multiplier follows a shifted-lognormal law $\xi_i = \min\!\bigl(\underline\xi + Y_i,\ \overline\xi\bigr)$, where $Y_i$ is lognormal with a base coefficient of variation of $0.15$, and the fixed bounds are $\underline\xi = V^{\mathrm{nom}}/90 \approx 0.889$ (speed limits) and $\overline\xi = V^{\mathrm{nom}}/50 = 1.6$ (a leg-average congestion floor). These same bounds $[\underline\xi,\overline\xi]$ define the box and budgeted uncertainty sets used by the robust methods of Sections~\ref{sec:ro}--\ref{sec:robu}. The realized energy on each leg follows from the ECR curve at the implied speed. \NEW{The location parameter of $Y_i$ is calibrated numerically so that $\mathbb E[\xi_i] = 1$ exactly after truncation, which keeps the nominal plan an unbiased benchmark; at the base coefficient of variation the upper cap absorbs about 1\,\% of the mass. Multipliers are drawn independently across legs.}

\begin{atodo}
\textbf{Independence across legs is the assumption most exposed to attack.} Legs
are 25\,km and a real congestion event spans several consecutive ones.
\texttt{settings.py} already implements an AR(1) driver
(\texttt{TRAVEL\_TIME\_AR1\_RHO}, currently 0). Either run $\rho = 0.4$ on the
reduced set --- correlation should hurt exactly the methods that assume
independence, namely ROBU's budget calibration and 2SP's scenario set, and leave
LA comparatively unharmed, which would itself be a result --- or state the
limitation and predict the direction in Section~\ref{sec:discussion}. Silence is
the one option a referee will not accept.
\end{atodo}

\subsection{Benchmark instances} \label{sec:instances}

A set of benchmark instances of increasing complexity is generated along a single long-haul corridor. Charging station stops are placed at regular average intervals of 60\,km, customer stops are randomly place around clusters centers inserted along the corridor. The instances are parameterized based on route distance, customer count, and time windows width as summarized in Table~\ref{tab:instances}. For each instance, service windows are generated from the deterministic MILP solved with average travel time: each window is centered on the service start of the found optimal solution, with half-width drawn independently for each customer from the class range in Table \ref{tab:instances}\FIX{, and the centre jittered by $\pm10\,\%$ of that half-width}. Instances whose deterministic model is infeasible under nominal data are discarded and regenerated. Each Route/Customers/Time Windows class contains 50 independent and randomly generated instances. The number of stops \DEL{ranges from 51 on average for the short route instances, to 168 for the long route instances}\FIX{averages 50 on the short routes, 96 on the medium and 167 on the long ones, of which 16, 32 and 58 are charging stations and 26, 56 and 101 are rest-area nodes}.

\begin{acmt}
Counts recomputed over all 1\,800 generated instance files. Rest-area nodes are
the \emph{majority} of stops (52--60\,\%): they carry no service, no charging and
no window, but they carry the full HoS state chain, and they are the main driver
of model size. That is currently invisible in the text, and it is also what makes
the $\Gamma$ values of Section~\ref{sec:robu} as large as they are.

\medskip
The $\pm10\,\%$ centre jitter matters too: without it every window would be
centred on a schedule the deterministic model can hit exactly, and a referee
would suspect the design favours the offline methods.
\end{acmt}

\begin{table}[ht]
\centering
\caption{Experimental design. Full factorial,
         $3 \times 3 \times 4 \times 50 = 1\,800$ instances.}
\label{tab:instances}
\small
\begin{tabular}{lll}
\toprule
\textbf{Factor} & \textbf{Levels} & \textbf{Values} \\
\midrule
Route class   & Short  & 800--1\,200\,km     \\
              & Medium & 1\,500--2\,500\,km  \\
              & Long   & 3\,000--4\,000\,km  \\
\midrule
Customers     & Few    & 1--3    \\
              & Medium & 4--6    \\
              & Many   & 7--15   \\
\midrule
Time window   & None   & no window         \\
 (half-width  & Tight  & $U[0.5, 1.0]$\,h  \\
 per customer)& Medium & $U[1.0, 3.0]$\,h  \\
              & Large  & $U[3.0, 6.0]$\,h  \\
\midrule
\NEW{Derived}   & \NEW{Stops $|\mathcal I|$}    & \NEW{50 / 96 / 167 (short / medium / long)} \\
              & \NEW{Budget $\Gamma$}         & \NEW{18 / 24 / 31} \\
\bottomrule
\end{tabular}
\end{table}

Each route instance class targets a distinct operational regime: \DEL{Small}\FIX{Short} instances complete with a single rest, Medium instances span two to three shifts, activating the rest, reduced-rest, and extension mechanisms. \DEL{Large instances approach the weekly driving cap.}\FIX{Long instances span four to five duty cycles, at which point the weekly provisions outside the scope of this model would begin to bind.}

\begin{acmt}
Two things: Small/Medium/Large here versus Short/Medium/Long everywhere else, and
a reference to a weekly cap the model does not enforce. Note also that ``Medium''
is simultaneously a route class, a customer class and a window class, so
``Medium/Medium/Medium'' is a real cell in Table~\ref{tab:gap}. Renaming the
customer classes C1--3 / C4--6 / C7--15 would remove the collision.
\end{acmt}

\subsection{Norwegian case study}
{\color{red} TBD }
In addition to the synthetic corridors, we evaluate all methods on a real long-haul instance along \prov{XXX} (\prov{XXX} km), with \prov{XXX} customer terminals at \prov{XXX} cities. Charger locations and rated powers are taken from already accessible charging stations for heavy duty vehicles, while leg distances and nominal speeds are obtained through Google maps API. Travel-time uncertainty on each leg is calibrated from \prov{XXX} (Statens vegvesen traffic data?)

\begin{atodo}
\textbf{Decision required now, not at submission.}
\texttt{src/instance\_gen/norway\_instance.py} is a loader skeleton: the code
path is complete but none of the four data files it needs exists
(\texttt{chargers.csv} from NOBIL, \texttt{customers.csv}, \texttt{legs.csv} from
OSRM, \texttt{dispersion.json} from Statens vegvesen). No Norwegian instance has
been generated or solved.

\medskip
\emph{Build it}: export HDV-capable sites from NOBIL for the E6 corridor, pick
3--5 terminal cities, pull leg distances from OSRM, calibrate the coefficient of
variation --- and ideally the AR(1) coefficient --- from
\texttt{trafikkdata.atlas.vegvesen.no}. That last item would simultaneously
discharge the independence objection above, which makes this worth more than a
case study normally is. The loader's own comment warns that corridor gradients
make the constant-ECR assumption optimistic here.

\medskip
\emph{Or cut it}, and remove ``and a real Norwegian long-haul instance'' from the
abstract. 1\,800 synthetic instances on a full factorial is a sufficient
computational study; an empty section is not. Either way the abstract must stop
promising it --- the version proposed on page 1 no longer does.
\end{atodo}
